\section{Overview}

KeyQuorum Lab runs the KeyQuorum crate in a browser as a small sandboxed
machine, compiled to WebAssembly. It has its own files, a SQLite store per
person, mock USB drives, and a mailbox relay answering in the page. Every
button and every terminal line runs the real \code{keyquorum} and
\code{keyquorum-device} commands --- there is no separate ``demo mode''
standing in for them.

\begin{noteBox}
This manual documents the lab's own GUI: its header, tabs, and guided
tutorials. It does not re-document \code{keyquorum} / \code{keyquorum-device}
themselves --- that is the separate \emph{KeyQuorum} manual, which the lab
runs underneath its interface unchanged. Where a control here is a thin
wrapper over a CLI concept (custody, quorum, private bridges, authenticated
updates, the mailbox relay), this manual names that concept without
re-explaining it.
\end{noteBox}

\Screenshot{docs/latex/images/lab/header.png}{The lab header: build commit,
the Tutorials \& Documentation and Reset Lab buttons, the disclaimer, and
the Active user bar.}

\textbf{Reset Lab} discards everything the session has done --- every
drive move, sent letter, provisioned key, password file, export, share
link, signature, and activity-log entry --- and replaces it with a fresh
copy of the seeded demo state, instantly. Nothing here is persisted
outside the browser tab, so this is the only way back to a known starting
point; it is also what \textbf{Exit tutorial} and the tutorial picker
leave alone; a gated tutorial in progress does not survive a reset.

Tasks that take several commands in a plain shell take one here --- a file is
sent with \code{keyquorum send}, a letter is opened and answered with
\code{keyquorum inbox open}, and \code{keyquorum doctor} says what is
missing --- and the lab shows each command it ran. The older, longer
procedures (\code{deliver}, \code{file share|receive|ack}, \code{relay pull})
still work if typed in the Terminal tab, but no button uses them, and they are
not recommended (Section~\ref{sec:lab-terminal}).

The lab itself is a browser-based virtual machine: its mock USB drives
emulate USB-token custody entirely in memory, and every slot passphrase is
a published default value (Section~\ref{sec:lab-default-credentials}), so
they give none of the physical protection real hardware keys do. What is
real is the software underneath: every action, in every tab, is one of the
same commands documented in the KeyQuorum manual.

\subsection{Default credentials}
\label{sec:lab-default-credentials}

Every credential the lab uses is a published placeholder, not a secret.
The lab ships as a public WebAssembly bundle that runs entirely in your
browser, so anything baked into it --- the seeded people, their USB
passphrases, and the seeded files --- is visible to anyone who inspects the
bundle. That is by design: none of it guards real data, so nothing here
needs to stay confidential, and no real file or hardware key is ever at
risk of passing through the lab.

\textbf{USB slot passphrases.} Every seeded USB slot unlocks with
\code{lab-demo-<label>}, where \code{<label>} is that slot's dotted tree
label:

\begin{center}
\begin{tabular}{@{}lll@{}}
\toprule
\textbf{Person} & \textbf{Label} & \textbf{Passphrase} \\
\midrule
Morgan & \code{M}     & \code{lab-demo-M} \\
Sarah  & \code{M.S}   & \code{lab-demo-M.S} \\
Alice  & \code{M.S.1} & \code{lab-demo-M.S.1} \\
Bob    & \code{M.S.2} & \code{lab-demo-M.S.2} \\
David  & \code{M.A}   & \code{lab-demo-M.A} \\
Emma   & \code{M.A.1} & \code{lab-demo-M.A.1} \\
Chris  & \code{M.A.2} & \code{lab-demo-M.A.2} \\
\bottomrule
\end{tabular}
\end{center}

\textbf{Terminal default values.} Typing a command yourself at the Terminal
tab --- rather than clicking a GUI action that stages your own value first
--- is answered with the same published defaults a real prompt would use,
whenever it asks for a password or PIN:

\begin{itemize}
  \item Password: \code{lab-demo-password}
  \item PIN: \code{0000}
\end{itemize}

\begin{noteBox}
Seeded quorum-protected files never use a password. They are protected by
KeyQuorum's quorum and device-custody policies instead, and are unlocked
with the USB identities above.
\end{noteBox}

\subsection{Identity and drives}

The \textbf{Active user} bar shows who you are acting as right now, and
which mock USB drive their personal key slot lives on. Clicking another
name's chip switches to them, the way sitting down at a different desk
would. Most actions need that person's drive \emph{inserted} first --- a
mock drive behaves like a real one: its slots are only usable while it is
connected, and moving a slot onto someone else's drive needs both drives
inserted at once, matching the physical requirement of moving a token
between two real USB drives.
